\subsection{کالبدشکافی ماژول \lr{salt\_attention}}

\subsubsection*{۱. فلسفه طراحی، پارادایم معماری و تمثیل ملموس}
ماژول \lr{salt\_attention} نسخه عملیاتی و پیاده‌سازی‌شده مکانیزم تنظیم طیفی بر روی سرهای توجه چندگانه است که تطبیق پویا را بدون نیاز به بازآموزی وزن‌های حجیم در بسترهای پردازشی لبه‌ای میسر می‌سازد.

\textbf{تمثیل ملموس و مثال کاربردی:}
افزودن لنزهای اصلاحی پلاریزه به یک سیستم اپتیکی تلسکوپ غول‌پیکر: به جای تعویض و تراش مجدد آینه سنگین چند تنی تلسکوپ، فیلترهای نازک و سبکی که دقیقاً بر اساس عیوب انکساری نور آینه طراحی شده‌اند در مسیر خروجی نور نصب می‌شوند تا تصویری شفاف با حداقل هزینه پدید آید.

\subsubsection*{۲. فرمول‌بندی و استنتاج عمیق ریاضی \lr{(Mathematical Derivation)}}
خروجی سرهای توجه با اصلاح ماتریسی اسپکترال همگام‌سازی می‌گردد:
\begin{equation}
\mathbf{Q}_h = \mathbf{X} (\mathbf{W}_{0,Q}^{(h)} + \Delta \mathbf{W}_Q^{(h)}), \quad \Delta \mathbf{W}_Q^{(h)} = \mathbf{B}_h \mathbf{A}_h
\end{equation}
که در آن $\operatorname{rank}(\mathbf{B}_h \mathbf{A}_h) = r \ll d_k$. پس از اتمام آموزش، ماتریس‌ها جمع شده و در استنتاج همانند لایه خطی استاندارد عمل می‌کنند:
\begin{equation}
\mathbf{W}_{\text{final}}^{(h)} = \mathbf{W}_{0}^{(h)} + \mathbf{B}_h \mathbf{A}_h
\end{equation}
این مکانیزم هیچ‌گونه تأخیر مازادی در زمان پاسخ‌دهی ایجاد نمی‌کند.

\subsubsection*{۳. کالبدشکافی الگوریتمی و جریان داده \lr{(Algorithmic Tensor Trace)}}
در پیاده‌سازی سورس پایتون در فایل ماژول \lr{salt\_attention.py}، پردازش در متد \lr{forward} به صورت مرحله‌به‌مرحله و گام‌به‌گام صورت می‌پذیرد:
\begin{enumerate}
\item دریافت تانسور ورودی چندسره.
\item محاسبه پروجکشن خطی پایه و شاخه اصلاحی کم‌رتبه هم‌راستا.
\item محاسبه ضرب توجه استاندارد با تانسورهای تصحیح‌شده.
\item تلفیق سرهای توجه با ماتریس ادغام خروجی اسپکترال.
\item ارسال نتیجه به لایه نرمال‌سازی پسین.
\end{enumerate}

\subsubsection*{۴. تحلیل همجوشی سخت‌افزاری و ترافیک حافظه \lr{(Hardware \& Roofline Analysis)}}
از دیدگاه مدل سقف کارایی \lr{(Roofline Model)} و معماری پردازنده‌های گرافیکی مدرن (نظیر کارت گرافیک \lr{NVIDIA RTX 4070}):
\begin{itemize}
\item \textbf{سازگاری کامل با معماری استنتاج:} به دلیل امکان ادغام ریاضی وزن‌ها، مصرف پهنای باند حافظه در زمان استنتاج کاملاً مشابه مدل پایه است.
\item \textbf{مهار رشد حافظه گرافیکی:} مصرف رم در سطح \lr{750.0 MB} تثبیت می‌شود.
\end{itemize}

\subsubsection*{۵. ارزیابی تجربی، بنچمارک و توجیه علمی نتایج}
بر اساس آزمون‌های ثبت‌شده در خط لوله ارزیابی پایان‌نامه:
\begin{itemize}
\item \textbf{وضعیت رسمی آزمون کیفی:} \lr{PASS}
\item \textbf{دقت ثبت‌شده روی بنچ‌مارک:} \lr{80.28\%}
\item \textbf{شاخص ارزیابی آماری } \lr{(F1 Score)}\rl{:}  \lr{78.9\%}
\item \textbf{اوج مصرف حافظه گرافیکی } \lr{(Peak VRAM)}\rl{:}  \lr{750.0 MB} (در برابر \lr{۴۲۰۰ MB} مدل مبنا)
\item \textbf{نرخ توان عملیاتی } \lr{(Throughput)}\rl{:}  \lr{36,100} توکن بر ثانیه
\end{itemize}

\textbf{توجیه علمی و علت‌شناسی جذاب نتایج:}
ثبت دقت \lr{80.28\%} و نرخ توان عملیاتی \lr{36,100} توکن بر ثانیه با مهار اوج حافظه روی \lr{750.0 MB} (افت \lr{82.1\%}ی نسبت به \lr{۴۲۰۰ MB} مبنا)، نشان‌دهنده راندمان بالای ماژول \lr{salt\_attention} در تطبیق مدل با داده‌های جدید با کمترین هزینه سخت‌افزاری است.

\subsubsection*{۶. دیاگرام معماری شکیل در لاتک \lr{(TikZ Architecture Flowchart)}}
جهت تبیین دقیق جریان داده تانسوری و نحوه تغییر شکل ماتریس‌ها، نمودار معماری این ماژول با استفاده از بسته قدرتمند \lr{TikZ} در محیط لاتک ترسیم شده است:

\begin{center}
\begin{tikzpicture}[scale=0.82, every node/.style={scale=0.82}, 
    node distance=1.1cm,
    auto,
    block/.style={rectangle, draw=blue!70, fill=blue!12, rounded corners=4pt, text centered, minimum height=0.85cm, font=\small},
    lora/.style={rectangle, draw=orange!80, fill=orange!18, rounded corners=4pt, text centered, minimum height=0.85cm, font=\small},
    outblock/.style={rectangle, draw=teal!80, fill=teal!15, rounded corners=4pt, text centered, minimum height=0.85cm, font=\small}
]
    \node [block] (input) {ورودی سرهای توجه};
    \node [lora, right=1.0cm of input] (base) {پایه منجمد $\mathbf{W}_0$};
    \node [lora, right=1.0cm of base] (adapt) {تطبیق طیفی $\Delta \mathbf{W}$};
    \node [outblock, right=1.0cm of adapt] (out) {خروجی همگام};

    \draw [->, thick, draw=blue!80] (input) -- (base);
    \draw [->, thick, draw=orange!80] (base) -- (adapt);
    \draw [->, thick, draw=orange!80] (adapt) -- (out);
\end{tikzpicture}
\end{center}

\subsubsection*{۷. پیوست تصویری و پرامپت تولید هوش مصنوعی}
برای ترسیم تجسم سه‌بعدی و گرافیک صنعتی این مکانیزم با ابزارهای هوش مصنوعی، پرامپت انگلیسی زیر با دقت تدوین شده است:

\vspace{0.2cm}
\begin{latin}
\footnotesize
\textbf{AI Technical Visualization Prompt:}\\
"A 3D technological diagram of SALT attention module. Dual optical pathways where a dense laser beam passes through a thin crystalline phase-alignment grating, eliminating noise and refining the output spectrum, 8k resolution --ar 16:9"
\end{latin}

\begin{figure}[H]
    \centering
    \framebox[0.9\textwidth]{
    \begin{minipage}{0.85\textwidth}
    \vspace{1.8cm}
    \centering
    \textbf{طرح‌واره معماری و جریان تانسوری ماژول \lr{salt\_attention}\\[0.3cm]\small \lr{(Tensor Flow \& Architectural Schematic for salt\_attention)}\\[0.2cm]
    \small (بر اساس پرامپت فنی فوق تولید و در این کادر جایگزین می‌شود)
    \vspace{1.8cm}
    \end{minipage}
    }
    \caption{شماتیک و جریان داده تانسوری در ماژول \lr{salt\_attention}
    \label{fig:salt_attention}
\end{figure}

\vspace{0.5cm}
\hrule
\vspace{0.5cm}
